% !PRINTED 608
\weberpage{608}

%<<CONT
\weq{4}{P = \prod\nolimits \frac{1}{\left[1-N(\mathfrak{p}_1)^{-s}\right]},}
so ist nach (2) und (3) für $s=1$
\weq{5}{(s-1)\,P^n \text{ endlich und von Null verschieden,}}
und nach § 163, 6.:
\weq{6}{(s-1)\,P^h \text{ endlich,}}
und wenn man hieraus $P$ eliminiert, so folgt:
\[
(s-1)^{n-h} \text{ endlich.}
\]
Also ist
\weq{7}{n \overline{\overline{>}} h.}
Daraus folgt der Satz:

8. \gesperrt{Der Relativgrad $n$ des Klassenkörpers kann nicht kleiner sein als die Klassenzahl $h$ der Gruppe. Ist $n=h$, so ist $(s-1)P^h$ für $s=1$ nicht Null, und die Summen $Q_2, Q_3, \ldots, Q_h$ sind gleichfalls von Null verschieden.}

Kann man eine Gleichung in $\Omega$ vom Grade $h$ bilden, aus der sich ein Körper ableiten läßt, der die charakteristischen Eigenschaften des Klassenkörpers erkennen läßt, und ist der Grad dieses Körpers nicht größer als $h$, so folgt also, daß er auch nicht kleiner sein kann, also $= h$ sein muß. Damit ist die Irreduzibilität der fraglichen Gleichung erwiesen.

Hierauf beruht der Nachweis der \gesperrt{Irreduzibilität der Kreisteilungsgleichung}, den wir im § 198 des II. Bandes gegeben haben. Dort trat an Stelle von $\Omega$ der Körper der rationalen Zahlen. Im nächsten Abschnitt werden wir noch weitere Anwendungen hiervon machen.

Ob es Gruppen gibt, bei denen der Grad des Klassenkörpers größer als die Klassenzahl ist, bleibt dahingestellt. Ich halte es für unwahrscheinlich.

Wir beweisen noch, daß es höchstens \gesperrt{einen} Klassenkörper geben kann. Jeder solche Körper wird durch Adjunktion einer algebraischen Zahl $\xi$ zu $\Omega$ erzeugt.

Es sei $\xi$ eine den Klassenkörper $\mathfrak{K} = \Omega(\xi)$ erzeugende \gesperrt{ganze Zahl}, die also einer in $\Omega$ irreduzibelen Gleichung
\weq{8}{f(\xi) = \xi^n + \alpha_1 \xi^{n-1} + \alpha_2 \xi^{n-2} + \cdots + \alpha_n = 0}
genügt, in der die $\alpha_1, \alpha_2, \ldots, \alpha_n$ ganze Zahlen in $\Omega$ sind.
